\section{Exact remaining gaps}\label{sec:22.2}

The decomposition theorem of \cref{sec:8.6} holds for a finite insulator. For the infinite crystal one estimate is missing, a superadiabatic generator built from exponentially localized terms, and every sector of every rate in this book is added on the strength of the finite-crystal statement with its stated error. At the other end is the silicon rate of a particle with an anapole moment. At \SI{10}{MeV} it exceeds the static-electron formula at least 16 times, and the local-field screening of the transverse current, which would turn that bound into a number, is open (\cref{sec:10.4}). The first statement needs a proof and changes no number; the second needs a calculation and changes a limit. Both have the same form, a strongest result and one step beyond it. \Cref{tab:22-gaps} collects all 74 such items by chapter, with the result each one limits.

\subsection{Exact remaining gaps and nuisances}\label{sec:22.2-definition}

\begin{keyresult}[Exact remaining gap]
An \emph{exact remaining gap} is a pair. Its first part is the strongest statement that the book derives about a quantity or a claim, with the conditions under which it holds: a theorem under stated hypotheses, a bound, an estimate with its factor, or a computed number that awaits an independent check. Its second part is the single missing step, a proof, a calculation or a measurement, that would turn the statement into the full result. Each gap is stated once, in the section where the quantity arises, which gives both parts and, for a calculation, points to its specification in \crefrange{chap:30}{chap:33}.
\end{keyresult}

A gap differs from a nuisance. A nuisance parameter of \cref{sec:8.8} carries a truncation whose next order is computed or estimated from a small parameter, and it enters a likelihood with a width; the attraction between electron and hole, which raises the silicon electron-recoil rate by $\delta_{\mathcal{X}} \approx +0.28$ over the \SI{3}{MeV} window, is one (\cref{tab:10-nuisances}). \Cref{sec:8.8} also names a third case, a truncation with neither a computed next order nor a small parameter, which receives no width. Every such truncation is an exact gap, and so is every claim the book can derive only under a hypothesis it cannot yet remove. The nuisance tables at the end of each chapter carry the first kind of item, and \cref{tab:22-gaps} the second. One item can have both a sized part and an open part. The attraction is such an item: its model size is a nuisance, and the finite-momentum Bethe--Salpeter calculation that would replace the model is a gap. The screening of the phonon-assisted sideband splits the same way.

What a gap limits decides how we can use it, and the table sorts the 74 into five classes. A gap that limits only a proof leaves every number of the book standing. Its missing step extends a derived statement beyond its conditions, from a finite crystal to an infinite one or from perturbation theory to all orders. Three gaps are of this kind, those of the decomposition theorem. The other 71 involve a number. For 10 of them the book gives a bound or a two-sided bracket, such as the anapole bound above. For 17 it gives an estimate, a model value or a factor that is not a bound, such as the phonon-assisted rate of silicon below the zero-phonon edge (\cref{sec:13.9}). For 39 the number is unsized: the book can say what the result depends on and how to compute or measure it, but assigns it no width. The last 5 are tests. The book computes the number, and the missing step is an independent check of it, a second set of force constants for the phonon-assisted onset or a cold measurement of the silicon absorption edge.

The missing steps are of three kinds. In 56 rows the step is a calculation specified in the section or in \crefrange{chap:30}{chap:33}, in 11 a measurement, and in 7 a proof or a method that does not yet exist, such as a controlled treatment of Dirac targets at $\tilde\alpha \gtrsim 1$ (\cref{sec:16.1}) or a real-time method for quantum nuclei at low temperature (\cref{sec:11.5}). Where a section names more than one step, the row gives the decisive one. \Cref{sec:22.1} orders the calculations among these steps by the results they would change and states the check each must pass.

\subsection{Gaps by chapter}\label{sec:22.2-table}

\begin{landscape}
{\scriptsize\setlength{\tabcolsep}{3pt}\setlength{\LTcapwidth}{\linewidth}
\begin{longtable}{@{}>{\raggedright\arraybackslash}p{0.16\linewidth}>{\raggedright\arraybackslash}p{0.38\linewidth}>{\raggedright\arraybackslash}p{0.36\linewidth}>{\raggedright\arraybackslash}p{0.07\linewidth}@{}}
  \caption{Exact remaining gaps of the book by chapter, each indexed to the one section that states it. Strongest statement: what that section derives, with its conditions. Missing step: the decisive proof, method, calculation or measurement that would complete it, with the section that specifies a calculation. Limits: \emph{proof}, no number of the book changes. \emph{Bounded}, a number is missing and the section gives a bound or a two-sided bracket. \emph{Estimate}, the section gives an estimate that is not a bound. \emph{Unsized}, the section assigns no size. \emph{Test}, the section computes the number and the missing step checks it. The statements are abbreviated; the section named in the first column gives each in full. Survivors: the eleven crystals with a documented detector record (\cref{sec:21.7}). The calculations that would change a printed result, limit or ranking are ordered in \cref{tab:22-calc}.}
  \label{tab:22-gaps}\\
  \toprule
  Gap; section & Strongest statement & Missing step & Limits \\
  \midrule
  \endfirsthead
  \toprule
  Gap; section & Strongest statement & Missing step & Limits \\
  \midrule
  \endhead
  \bottomrule
  \endfoot
  \multicolumn{4}{@{}l}{\textit{\Cref{chap:8}. The response framework}} \\
  \midrule
  Sector labels outside the off-shell window; \cref{sec:8.6} & Split into zero-phonon and phonon-dressed sectors, and convention independence of the totals, derived in perturbation theory in $g$ with distributional limits; Dyson convention \cref{eq:8.6-dyson} & A limiting-absorption principle for the coupled electron--phonon Hamiltonian at fixed total crystal momentum, uniform in the coupling and the number of cells; it would also sharpen \cref{eq:8.6-bin} from $(\varGamma\Delta E_{\mathrm{c}})^{1/2}$ to $\varGamma$ & proof \\
  Uniformity in the number of nuclei; \cref{sec:8.6} & Superadiabatic leak \cref{eq:8.6-leak} of order $\kappa_{\mathrm{BO}}^{2(n+1)}\alpha^2m_e$, proved for a finite crystal & A superadiabatic generator built from exponentially localized terms, so that the leak per unit volume stays finite & proof \\
  Dressing of gapless targets; \cref{sec:8.6} & Bin bound \cref{eq:8.6-bin} for any orthogonal sectors given a dressed vacuum; fails for $U_{\mathrm{qp}} = 1$; connected part of $\mathfrak{h}_{\mathrm{off}}$ estimated at \SI{0.2}{eV} in aluminum, with no commutator bound uniform in the number of cells & Existence of $U_{\mathrm{qp}}$ beyond perturbation theory, with bounds & proof \\
  \midrule
  \multicolumn{4}{@{}l}{\textit{\Cref{chap:9}. Channel atlas}} \\
  \midrule
  Symmetry-allowed cells of the channel atlas; \cref{sec:9.1} & Eight cells of \cref{tab:9-atlas} allowed at leading order and computed in no section: a spin coupling to core electrons, Cooper pairs, Migdal transitions, Dirac pairs, phonon-assisted pairs, dopant transitions and Frenkel excitons, and the electron density acting on a freely recoiling atom through its following cloud, whose rigid-atom density $f_\kappa(q)$ is used in the multiphonon vertex (\cref{sec:11.3}) & The eight rates, each from the response of the home section of its row with the vertex of its column & unsized \\
  \midrule
  \multicolumn{4}{@{}l}{\textit{\Cref{chap:10}. Electronic excitations}} \\
  \midrule
  Local fields of the transverse current; \cref{sec:10.4} & Static-electron formula low by at least 16 at \SI{10}{MeV} for independent particles, at least 29 with the density screened (\cref{eq:10.4-band}); current screening 4--\SI{15}{\percent} in a model, 10--\SI{30}{\percent} estimated, both lowering & Mixed-block screening of \cref{eq:D.30-dyson}, which fixes the in-medium ratio $r_{\mathrm{med}}$ (\cref{app:E.16}) & bounded \\
  Finite-$q$ Bethe--Salpeter response with coupling block and bound lines; \cref{sec:10.5} & Attraction $\approx +0.28$ over the \SI{3}{MeV} window from a model, coupling block about \SI{10}{\percent} (\cref{eq:10.5-coupling}), resonant-only evaluation off by two (\cref{eq:10.5-res}), lines from \cref{eq:10.5-lines} & Coupled calculation for Si and Ge with the lines added from computed masses and vertices, $f$-sum to a percent (\cref{app:E.15}) & estimate \\
  Triplet beyond independent particles; \cref{sec:10.5-triplet} & Same lines and selection rules as the singlet, no local fields, no screening; 2.2 times the independent-particle spectrum below \SI{3}{eV} in a model & Ladder with the spin-channel exchange--correlation kernel; control: a quarter of the independent density response without interaction (\cref{app:E.16}) & estimate \\
  Self-trapped excitons of NaI; \cref{sec:10.5-open} & Poisson replica series with Huang--Rhys factor $S_{\mathrm{HR}}$ for linear coupling, a harmonic lattice and a fixed vertex (\cref{eq:E.15-ste}) & $S_{\mathrm{HR}}$, $\bar\omega$ and $\rho_{\lambda0}(\vect{q})$ at keV momenta (\cref{app:E.15}) & unsized \\
  Four-component finite-$q$ response; \cref{sec:10.6-open} & At most \SI{2}{\percent} for Ge $3d$, central value near $10^{-3}$; $10^{-6}$ beyond two components & A calculation with no-pair projection, the $3d$ binding as a nuisance, checked on the relativistic $f$-sum shift (\cref{app:E.15}) & bounded \\
  Spin--orbit mixing at the Ge edge; \cref{sec:10.6-open} & Computed and small in Si; up to \SI{10}{\percent} in amplitude near the Ge L minimum, estimated & Ge bands with spin--orbit coupling in the measured valley order (\cref{app:E.15}) & estimate \\
  Interference of contact and second-order amplitudes; \cref{sec:10.8-open} & Linear in the quadratic response, a three-point function with no reduction to two-point responses; known for the electron gas (Barkas term) & Quadratic response of the crystal, in the local-density mapping of \cref{app:E.26} or at second order in time-dependent perturbation theory & unsized \\
  Bound shells in dark Compton scattering; \cref{sec:10.8-open} & Second-order terms of \cref{eq:10.8-TC} vanish away from the shells; nonzero for Si K at 45 to about \SI{100}{keV}, Ge K from \SI{118}{keV}, L shells below about \SI{20}{keV} & Compton profiles of the bound shells at these energies, from the large-$q$ density response (\cref{app:E.26}); the refractive index enters as an input & unsized \\
  Fermionic absorption on bound shells; \cref{sec:10.9-open} & Free-electron vertex low by at least $1/r$, 1.7 for xenon $1s$ at \SI{56.5}{keV}; the transverse current only raises it; sign unknown only at order $(Z\alpha)^2 \approx 0.155$ & Transverse current response with Dirac--Hartree--Fock states, and the published limits redone (\cref{app:E.26}) & bounded \\
  Retarded density row of boosted dark matter; \cref{sec:10.9-open} & Heavy mediator: limit unchanged; massless mediator: shift below \SI{1}{\percent} for $\mchi \ge \SI{10}{keV}$; density-row ratio 0.05 to 0.93 at \SI{1}{keV}, undetermined below about \SI{3}{keV} (\cref{sec:21.5-corrected}) & The published calculation redone with the exact density row and the transverse row, Cherenkov pole included (\cref{app:E.16}) & bounded \\
  \midrule
  \multicolumn{4}{@{}l}{\textit{\Cref{chap:11}. Lattice excitations}} \\
  \midrule
  Valence deformation above the zone; \cref{sec:11.3} & Translation sum rule on the cell sum; rigid-atom vertex $Z_\kappa - f_\kappa(q)$; core share computed (\cref{app:E.17}); nucleon couplings unaffected & All-electron density response of displaced atoms in a large cell, 3 to \SI{20}{keV}, tested on the translation sum rule (\cref{app:E.17}) & unsized \\
  Strongly anharmonic targets; \cref{sec:11.5} & Exact moments by path integrals bound the cumulative weight (Markov--Krein band); Kubo-transformed dynamics only for $T \gtrsim \omega/10$; neutron Compton scattering would measure the momentum distribution & A real-time method for quantum nuclei, exact at low temperature (\cref{app:E.19}) & bounded \\
  Cold in-band dielectric function of Si; \cref{sec:11.6} & $\alpha(\omega; T \to 0) \le \alpha(\omega; \SI{293}{K})$ (\cref{eq:11.6-Tbound}): isotope channel at least $5.0\times10^3$ times the electromagnetic one at \SI{50}{meV}; computable from the derivatives of the Born charges (\cref{app:E.17}) & Two-phonon absorption below \SI{65}{meV} measured at low temperature & bounded \\
  Cubic harmonics of the Si loss function; \cref{sec:11.7} & Pseudopotential-model components to $2.8\,q_0$; local fields change the direction average by 1.3 to 1.5 at \SI{10}{MeV} & First-principles loss function with local fields on Lebedev grids of increasing order (\cref{app:E.16}) & estimate \\
  \midrule
  \multicolumn{4}{@{}l}{\textit{\Cref{chap:12}. The Migdal effect}} \\
  \midrule
  Non-rigid Wannier ratio; \cref{sec:12.3} & Wannier following exact for core shells and flat bands, bounded only in absolute terms for the valence band of Si & Migdal probability with the self-consistent vertex against rigid following, checked on the acoustic and drag-mass sum rules (\cref{app:E.18}) & bounded \\
  Off-diagonal dielectric elements; \cref{sec:12.3-validity} & Scalar dielectric function high by about \SI{35}{\percent}; $\vect{G} \neq \vect{G}'$ elements raise the probability by 1.38 in a model, about 1.2 extrapolated, the lower end not a bound & Converged calculation with the full dielectric matrix (\cref{app:E.18}) & estimate \\
  Contact density at finite $q$; \cref{sec:12.5} & Known for spherical atoms; valence anisotropy $1.3\times10^{-3}$ in Si; relativistic normalization $\mathcal{R}_{\mathrm{rel}}$ (\cref{eq:12.5-Rrel}) with a quadrupolar part changed by a few percent & All-electron Dirac calculation of the crystal with a finite nucleus (\cref{app:E.17}) & estimate \\
  Site-resolved dynamical Born charges; \cref{sec:12.6-open} & Sum fixed by \cref{eq:12.6-sum}; GaAs bremsstrahlung carries the band of \cref{fig:12-brems} & Frequency-dependent charges per site (\cref{app:E.17}) & bounded \\
  Scalar absorption with couplings not proportional to mass; \cref{sec:12.6-open} & Needs $S^{ee}_{\vect{G}\vect{G}'}$ at vanishing reduced momentum & That response (\cref{app:E.18}) & unsized \\
  The Migdal chain as a whole; \cref{sec:12.6-open} & Chain from the kick to the electronic spectrum in place (\cref{sec:12.1,sec:12.5}); no crystal Migdal signal observed & Tagged-recoil measurement in Si: recoil energy and direction fixed by a scattered neutron, electrons counted & test \\
  \midrule
  \multicolumn{4}{@{}l}{\textit{\Cref{chap:13}. Phonon-assisted electron recoil}} \\
  \midrule
  Exciton--phonon replicas; \cref{sec:13.8} & Replica edges at $E_{\mathrm{min}}^{\mathrm{PA}}(\vect{q}) - E_x^b/n^2$, square-root rise at fixed $\vect{q}$; $1s$ transverse acoustic replica at \SI{1.1728}{eV}; completeness reduced to the two-pair sector & Bethe--Salpeter solution on Wannier-interpolated states in the gauge of the vertices, about $10^3$ points of $\vect{k}$ per pocket (\cref{app:E.15}) & estimate \\
  Onset on one set of force constants; \cref{sec:13.9-open} & \SI{0.3332}{MeV} (replica), \SI{0.3374}{MeV} (free pair), set by $\omega_{\mathrm{TA}}(\vect{Q}^*)$ and the gap & $\omega_{\mathrm{TA}}(\vect{Q}^*)$ at two lattice constants and with a second pseudopotential (\cref{app:E.14}) & test \\
  Zero-phonon to phonon-assisted ratio; \cref{sec:13.9-open} & Phonon-assisted lead to $\approx\SI{0.43}{MeV}$ (central), at least $\approx\SI{0.42}{MeV}$; two channels 0.43--\SI{0.51}{MeV}; zero-phonon lead above; both rates estimated & Ratio of the two rates on common bins, with the full phonon-assisted amplitude and the complete zero-phonon edge (mirror-plane cones, second heavy hole, X face) (\cref{app:E.14}) & estimate \\
  Rate and reach at first-principles level; \cref{sec:13.9-open} & Estimates normalized to the measured cold edge (\cref{tab:13-window}) & The same calculation with its controls: optical limit, symmetry zeros, soft acoustic limit, completeness (\cref{app:E.14}) & estimate \\
  Plane-wave $f$-sum check; \cref{sec:13.9-open} & 0.81 of $q^2/2m_e$ with 16 bands, 0.89 with 100, at 0.1--\SI{0.5}{\angstrom^{-1}} & Separation of the omitted bands from the nonlocal term (\cref{app:E.14}) & test \\
  Germanium window and phonons; \cref{sec:13.9-open} & Phonon-assisted onset 0.212--\SI{0.214}{MeV} from the measured transverse acoustic branch; zero-phonon edge not known, so the germanium maps of \cref{sec:17.1-ge-gaas} have open boundaries & Quasiparticle bands with spin--orbit coupling in the measured valley order (\cref{app:E.14}) & unsized \\
  Screening beyond the optical calibration; \cref{sec:13.9-open} & External factor and local-field wings 0.3--2.6 relative to the calibration, estimated; pair-channel dependence on $q$ unsized & Supercell calculation with local fields in the supercell basis (\cref{app:E.16}) & estimate \\
  GaAs control of the window; \cref{sec:13.9-open} & Phonon-assisted band 0.43 to 0.52--\SI{0.56}{MeV} ($E_{\mathrm{L}}(0) = \SI{1.815}{eV}$) in the parabolic description & Band calculation and rate with the polar vertex (\cref{app:E.14}) & estimate \\
  Cold absorption edge of Si; \cref{sec:13.9-open} & Untransferred \SI{77}{K} data off by more than two near the edge; two cold datasets differ by twenty at \SI{1.20}{eV} & Absorption within \SI{25}{meV} of threshold at millikelvin temperatures & test \\
  \midrule
  \multicolumn{4}{@{}l}{\textit{\Cref{chap:14}. Magnetic excitations}} \\
  \midrule
  Two-magnon interactions and bound pairs; \cref{sec:14.2} & Ladder \cref{eq:14.2-ladder} with the vertex from the spin Hamiltonian; interaction $1/2|\vect{S}| = 0.5$ in NiO, not small; existence of a bound pair below the continuum open & Solution of the ladder for NiO on $30^3$--$40^3$ wavevectors (\cref{app:E.20}) & unsized \\
  Longitudinal spin response from neutrons; \cref{sec:14.4} & Transverse block fixes the transverse couplings; for $\mathcal{O}_4$ it is a lower bound; for $\mathcal{O}_6$, $\mathcal{O}_{10}$ one $\vect{q}$ gives nothing; without the isotropy hypothesis of \cref{sec:14.4-iso}, broken only by spin--orbit and dipolar terms, only a model of the spin density & A neutron test of that hypothesis, turning the moments at fixed $\vect{q}$ & bounded \\
  Chirality after decay; \cref{sec:14.5} & GaAs lowest-band $\langle|\vect{L}^{\mathrm{ph}}|\rangle = 0.124$; below \SI{0.8}{MeV} that band takes under \SI{1}{\percent} of the rate & Survival of chirality through three-phonon decay, $(\mathbb{1} - \mathsf{M})^{-1}\vect{b}$ (\cref{app:E.19}) & unsized \\
  Large phonon magnetic moments; \cref{sec:14.5} & Born-charge moments at most $0.030\,\mu_N$ & Amplitudes of moments of order $\mu_B$ from the electronic orbital and spin coupling & unsized \\
  Optical activity at keV momenta; \cref{sec:14.5} & $q\,g_{33} = 0.034$ at \SI{1}{keV} in quartz, a scale only & Current response at $\pm\vect{q}$ with spin--orbit coupling and the nonlocal vertex, for tellurium or quartz (\cref{app:E.20}) & unsized \\
  Multiferroic response near ordering; \cref{sec:14.5} & Landau form within \SI{10}{\percent} of mean field only for $T \ge 0.94\,T_{\mathrm{CM}}$; spin waves controlled for $T \le 0.37\,T_{\mathrm{CM}}$ (\cref{lim:14.5b}) & Amplitude $\chi_\times$ and scaling function $\Phi$ of the dynamic-scaling form for a real material near $T_{\mathrm{CM}}$ & unsized \\
  \midrule
  \multicolumn{4}{@{}l}{\textit{\Cref{chap:15}. Superconductors}} \\
  \midrule
  Amorphous WSi; \cref{sec:15.3-wsi} & Rate needs the configurational average of the loss function, which differs from that of the averaged $\varepsilon$ at the order of its variance & Disorder-averaged loss function at keV momenta and the distribution of the local gap (\cref{app:E.20}) & unsized \\
  Strong-coupling Holstein factor; \cref{sec:15.4-strong} & Normal metal at $\omega \ll \omega_{\mathrm{D}}$: factor $(1 + \lambda_{\mathrm{tr}})^{-2} = 0.52$ in aluminum & Whether the factor survives in the superconducting form near the edge; rate with and without it (\cref{app:E.20}) & estimate \\
  Phonon-assisted spin current near the edge; \cref{sec:15.4-spin} & Coherence factor differs from \cref{eq:15.4-T}; competes with the direct term only above about \SI{10}{meV} & Its size near the edge & unsized \\
  Scalar absorption near the edge; \cref{sec:15.4} & The kinetic row couples to the amplitude channel of the pair field & Conserving response of \cref{sec:15.1} extended to that channel & unsized \\
  \midrule
  \multicolumn{4}{@{}l}{\textit{\Cref{chap:16}. Novel targets}} \\
  \midrule
  Dirac targets at strong coupling; \cref{sec:16.1} & Ladder controlled at first order in $\tilde\alpha$; BNQ-TTF has $\tilde\alpha \approx 10$ & A controlled method at $\tilde\alpha \gtrsim 1$ & unsized \\
  Depolarization of quantum dots; \cref{sec:16.3} & Factor $9|\varepsilon_{\mathrm{in}}|^2/|\varepsilon_{\mathrm{in}} + 2\varepsilon_{\mathrm{out}}|^2$ (\cref{eq:16.3-depol}), above one for real $\varepsilon_{\mathrm{in}} > \varepsilon_{\mathrm{out}}$; also computable (\cref{app:E.21}) & $\varepsilon_{\mathrm{in}}(\omega)$ of real dots over 2.5--\SI{10}{eV}, measured & unsized \\
  Impurity band; \cref{sec:16.5-band} & Isolated-donor response wrong at order one near the ionization edge for $\nd_{\mathrm{D}}/\nd_{\mathrm{c}} \approx 0.5$ & Disorder average on $10^4$--$10^5$ random sites with multivalley envelopes, checked on clusters (\cref{app:E.21}) & unsized \\
  Local-field ratio of a donor; \cref{sec:16.5} & Expected close to one in effective-mass theory & Dielectric matrix of the host and Bloch functions at the minimum (\cref{app:E.6}) & estimate \\
  Color-center yields at millikelvin; \cref{sec:16.6} & Yield $\langle N_{\mathrm{X}}\rangle p_{\mathrm{FH}}p_{\mathrm{surv}}f_{\mathrm{opt}}$; four factors open; \SI{5.9}{keV} X-rays and photodissociation at millikelvin would calibrate single factors & Tagged-neutron measurement on LiF at \SI{10}{mK}, counting centers per recoil & unsized \\
  Tunneling-model parameters of films; \cref{sec:16.7-open} & Loss tangent and acoustic loss fix $\bar Pp_{\mathrm{TLS}}^2$ and $\bar P\gamma^2$, not $\bar P$; $\gamma_{\mathrm{l}}/\gamma_{\mathrm{t}}$ 1.44--1.84 in 18 glasses; distributions for a model glass computable (\cref{app:E.22}) & One film measured three ways & unsized \\
  Interacting two-level systems; \cref{sec:16.8-open} & Spectral diffusion leaves a flat ensemble intact; branching gives mean cascade $1/(1 - b_{\mathrm{br}})$ and $N_{\mathrm{casc}}^{-3/2}$ as $b_{\mathrm{br}} \to 1$ & $b_{\mathrm{br}}$ from the barrier density and the ballistic phonon flux & unsized \\
  \midrule
  \multicolumn{4}{@{}l}{\textit{\Cref{chap:17}. Comparison across channels}} \\
  \midrule
  Coherent Migdal response at small $q$; \cref{sec:17.1} & Diagonal soft limit (\cref{sec:12.4}); its unsized factor acts on the whole Si Migdal band at 0.32--\SI{1}{MeV}, on a third of the rate at \SI{2}{MeV}, on at most \SI{7}{\percent} from \SI{3}{MeV} & Coherent response between about \SI{0.75}{keV} and $2q_{\mathrm{BZ}}$ & unsized \\
  Electron-following density at finite $q$; \cref{sec:17.1} & Long-wavelength model of \cref{sec:10.2}: acoustic peak and electron recoil equal at \SI{12.3}{MeV}, a property of the model & Following density to several zone momenta with Umklapp and multiphonon terms (\cref{app:E.16}) & unsized \\
  \midrule
  \multicolumn{4}{@{}l}{\textit{\Cref{chap:18}. Ionization yield}} \\
  \midrule
  Constant term and variance; \cref{sec:18.2-open} & $\bar n = \omega/\varepsilon_{\mathrm{eh}} + o(\omega)$ (\cref{eq:18.2-linear}); constants computed per seed & Second-order renewal estimates for the walk of \cref{eq:18.2-tagged} and the forgetting rate of the seed (\cref{we:22-constant}) & unsized \\
  Renewal estimate as a lower value; \cref{sec:18.2-open} & $F_{\mathrm{ren}} \le F_{\mathrm{eh}}$ in every computed case and in random parking & A proof for the class, or a counterexample & unsized \\
  Capture and terminal energies; \cref{sec:18.5-open} & Thinning law \cref{eq:18.5-thingen} given $f(x)$ & Computed capture probability $f(x)$, terminal-energy distribution and final pair separations (\cref{app:E.24}) & unsized \\
  Impact ionization beyond first order; \cref{sec:18.5-open} & A constant factor 0.8 moves $\varepsilon_{\mathrm{eh}}$ by about \SI{3}{\percent}, removed by calibration & An energy-dependent residue and exchange factor, which changes $F_{\mathrm{eh}}$ & unsized \\
  Memory of crystal momentum; \cref{sec:18.5-open} & Keeping probability below 0.1 within \SI{1}{eV} of threshold, one half above 3--\SI{5}{eV}; memory negligible in the deformation-potential model & The same with first-principles couplings & estimate \\
  Exciton and trapping rates at millikelvin; \cref{sec:18.5-open} & Binding energies and liquid-helium Auger lifetimes known; luminescent fractions and droplet thresholds at millikelvin also open & Binding, capture and trapping cross-sections of pairs measured at 10--\SI{50}{mK} (\cref{app:E.24,app:E.25}) & unsized \\
  \midrule
  \multicolumn{4}{@{}l}{\textit{\Cref{chap:19}. Low-energy nuclear recoils}} \\
  \midrule
  Charge yield below the binary threshold; \cref{sec:19.2-open} & No calculation predicts the Si yields and Fano factors there or the Ge point at \SI{254}{eV} & Close-collision kernel \cref{eq:19.2-kernel} for Si--Si and Ge--Ge with its exchange--correlation spread (\cref{app:E.23}) & unsized \\
  Release of close pairs at millikelvin; \cref{sec:19.6-open} & Energy-gap and instanton pictures both admitted for Si and Ge Frenkel pairs and Al split interstitials & Tunneling exponent along the computed path and $E_{\mathrm{reorg}}/\varepsilon^{\mathrm{hop}}$ to about \SI{5}{\percent} (\cref{app:E.22}) & unsized \\
  Stored energy in ionic crystals; \cref{sec:19.6-open} & Electronic route prompt; stored energy per deposit = excitons $\times$ conversion $\times$ survival $\times$ configurational energy of the pair & Conversion and close-pair survival probabilities at millikelvin & unsized \\
  Size law and rate of the excess; \cref{sec:19.6-lee} & Borel size law from branching; rate $\propto U_{\mathrm{rel}}$; spectral shape and time dependence factorize; $b_{\mathrm{br}}$ as for two-level systems & $U_{\mathrm{rel}}$ with its spread of rates for a holder and glue (\cref{app:E.25}) & unsized \\
  Carrier-driven defect reactions; \cref{sec:19.6-open} & Probability per capture \cref{eq:19.6-prxn}, a Poisson tail & $S_{\mathrm{HR,rxn}}$, $\omega_{\mathrm{rxn}}$ and the barrier of real defects with their charge states & unsized \\
  Sufficiency of the marks; \cref{sec:19.6-open} & Error of a photon-built detector response at most the total-variation distance at fixed $(N_{\mathrm{eh}}, \mathfrak{m})$, zero when the marks are sufficient & That distance for the features left out of $\mathfrak{m}$ (\cref{app:E.25}) & unsized \\
  \midrule
  \multicolumn{4}{@{}l}{\textit{\Cref{chap:20}. Charge, phonon and quasiparticle readout}} \\
  \midrule
  Trap emission and space charge; \cref{sec:20.1-open} & Emission rate = phonon flux, carrier flux and illumination times cross-sections; \SI{45}{meV} phonons decay within about \SI{1}{ns} & Phonon- and impact-ionization cross-sections of neutral and over-charged centers (\cref{app:E.25}) & unsized \\
  Luke spectrum beyond longitudinal phonons; \cref{sec:20.1-open} & Rate \cref{eq:20.1-luke} with longitudinal phonons, one deformation potential; drift velocities test only the momentum loss & Both acoustic branches with valley-resolved deformation potentials, checked against a phonon signal at fixed field & unsized \\
  Validation of transport; \cref{sec:20.1-open} & Hecht relation, charge sum, bias work \cref{eq:20.1-work}, Fano factor \cref{eq:20.1-Fpos} & An engine passing them to a few percent (\cref{app:E.25}) & test \\
  Thick-film offset of the quasiparticle yield; \cref{sec:20.4-open} & Cascade and kinetic equations differ by 0.04 above four gaps; $Y_{\mathrm{qp}}$ measured at fixed phonon energy would test it & Depth-resolved film transport at the measured $\alpha^2F$ & estimate \\
  Return of the gap energy at zero field; \cref{sec:20.5-routes} & Two-window contrast $0.28(1 - f_{\mathrm{ret}})$ at \SI{100}{eV}; sign changes at $t_1 = 3.87\,t_{\mathrm{rel}}$ for bimolecular against exponential release & Release-time law and site at millikelvin, from the rates of \cref{sec:18.5-open} (\cref{app:E.25}) & unsized \\
  \midrule
  \multicolumn{4}{@{}l}{\textit{\Cref{chap:21}. Exclusion limits and next targets}} \\
  \midrule
  Loss functions of the survivors; \cref{sec:21.7-open} & Migdal figure of merit between the plasmon-pole and gap-weighted limits & Random-phase response with local fields at the measured gap, tested on the inverse fourth moment (\cref{app:E.16}) & bounded \\
  Elastic constants; \cref{sec:21.7-open} & Computed tensors offset by \SI{4.9}{\percent} on average; TeO$_2$ beats CaWO$_4$ with probability 0.955 instead of 0.987 ; the other computed-tensor crystals and the cold density of CsI likewise open & Measured elastic constants of TeO$_2$ (\cref{app:E.17}) & estimate \\
  Internal backgrounds; \cref{sec:21.7-open} & Allowed or unique $\beta$ shapes with the relativistic Fermi function for $^{210}$Pb and $^{210}$Bi; first-forbidden $^{210}$Bi and $^{115}$In shapes not computed; $^{123}$Te bounded; cosmogenic rates measured for Si and Ge only (\cref{app:E.27}) & $^{210}$Pb activity of lead that can be grown into a crystal & estimate \\
  \midrule
  \multicolumn{4}{@{}l}{\textit{\Cref{chap:22}. Open problems}} \\
  \midrule
  Charged-current absorption; \cref{sec:22.3} & Allowed approximation: nuclear strengths times the Fermi function of the daughter (\cref{eq:21.7-fermi}); beyond it the rate is not a dark-matter factor times a target response & Shell-model Fermi and Gamow--Teller strengths per isotope, with quenched axial coupling, for every final state below $\mchi$ & unsized \\
  Charge-density-wave collective modes; \cref{sec:22.3} & Phason as the Goldstone pole and amplitudon of the conserving response on the mean field & That conserving response on a first-principles mean field with the self-consistent distortion & unsized \\
\end{longtable}}
\end{landscape}

The three proof gaps concern the sector decomposition on which every rate of the book is built. \Crefrange{chap:10}{chap:13}, on electronic, lattice, Migdal and phonon-assisted responses, hold 29 gaps, of which 6 are unsized: the response of a crystal at clamped or harmonic nuclei is a well-posed calculation, and where the book has not carried it out it bounds or estimates the result. \Crefrange{chap:18}{chap:20}, on the yield, the low-energy recoil and the readout, hold 17, of which 14 are unsized. What is missing there is mostly a rate at millikelvin temperatures, of capture, trapping, exciton binding, defect release or charge emission, for which neither a measurement nor a calculation exists, and the renewal analysis of the cascade, whose constants the book computes for each model but cannot size for the class.

Several gaps share their missing step. A finite-momentum Bethe--Salpeter calculation for silicon and germanium, with the coupling block and the bound lines, closes the attraction and triplet rows of \cref{sec:10.5}, and with exciton--phonon vertices in the same gauge also the replicas of \cref{sec:13.8}. Quasiparticle bands of germanium with spin--orbit coupling close the spin--orbit row of \cref{sec:10.6} and the germanium window of \cref{sec:13.9}, and with them the open boundaries of the germanium maps of \cref{sec:17.1}. The branching ratio $b_{\mathrm{br}}$ of strain-coupled defects is the missing input both of the avalanche question for two-level systems and of the size law of the low-energy excess (\cref{sec:16.8-open,sec:19.6-lee}). The millikelvin exciton and trapping rates of \cref{sec:18.5} also fix the release-time law of the gap energy in \cref{sec:20.5}. The 74 rows therefore stand for fewer independent steps. \Cref{tab:22-calc} lists once each calculation that would change a printed result, limit or ranking; the other calculation steps appear only here.

Eleven gaps need a measurement, and each section specifies it. Two of them test a chain that the book derives from the kick to the count: a tagged-recoil Migdal measurement in silicon, and the cold absorption edge of silicon within \SI{25}{meV} of its threshold. A third uses the same tagged-neutron method on F centers in LiF at \SI{10}{mK}, where counting single centers per tagged recoil measures the product of four yields as a function of the recoil energy (\cref{sec:16.6}). The other eight supply inputs that no calculation in the book replaces. On the side of the responses they are the cold two-phonon absorption of silicon in its one-phonon band, a neutron test of the isotropy of the spin response of a magnet, the dielectric function of real quantum dots and three measurements of one amorphous film. On the side of the detector and the crystal they are the millikelvin rates of excitons and traps, the stored-energy yields of ionic crystals, the elastic constants of TeO$_2$ and the $^{210}$Pb activity of lead that can be grown into a crystal.

\begin{workedexample}{The constant term of the mean pair number}\label{we:22-constant}
The constant term of the mean pair number is the first open item of \cref{sec:18.2-open}.

\emph{The strongest statement.} For any energy-conserving cascade whose share law becomes independent of the energy fast enough at high energy, is not concentrated on a lattice in $\ln x$ (non-arithmetic, in the language of renewal theory) and is regular near threshold, the mean pair number is $\bar n(\omega) = \omega/\varepsilon_{\mathrm{eh}} + o(\omega)$, \cref{eq:18.2-linear}, proved in \cref{app:D.44}. In the calibrated silicon cascade $\varepsilon_{\mathrm{eh}} = 3.232\,\Eg = 3.232\times\SI{1.160}{eV} = \SI{3.749}{eV}$.

\emph{What the book computes without the gap.} The cascade of \cref{sec:18.3} gives the constant $b_N$ in $\bar n = \omega/\varepsilon_{\mathrm{eh}} + b_N$ by direct computation for each seed: 0.35 for a seed that gives all the energy to one carrier, 0.008 for a uniform split of the energy. With the one-sided constants, $b_N = 0.35$ and $c_N = -0.02$ in the variance, the expansion matches the computed pair energy within \SI{0.07}{\percent} and the computed Fano factor within \SI{0.9}{\percent} at $\omega = 10\,\Eg$.

\emph{The missing step.} Second-order renewal estimates for the walk of $\ln x_j$ in \cref{eq:18.2-tagged}, together with the rate at which the cascade forgets its seed. They would show that $b_N$ exists for the whole class, as the average over the periodic remainder that a one-phonon model leaves, and that $\operatorname{Var}N_{\mathrm{eh}} = F_{\mathrm{eh}}\omega/\varepsilon_{\mathrm{eh}} + O(1)$.

\emph{What it limits.} At $\omega = \SI{10}{eV}$ the leading term is $10/3.749 = 2.667$ pairs. The one-sided seed adds 0.35, a mean of 3.017 and \SI{13.1}{\percent} of the leading term; the uniform split adds \SI{0.3}{\percent}, and the mean $\omega/\varepsilon_{\mathrm{eh}} + 0.19$ of the step model of \cref{sec:18.3} adds \SI{7.1}{\percent}. Each of these numbers comes from a computed cascade, not from the expansion, and none waits for the missing step. The one-sided constant falls below \SI{1}{\percent} of the leading term above \SI{131}{eV}, so the expansion matters only at a few pairs, where the cascade is cheap to compute. What the gap leaves without a size is the constant term of a cascade outside the computed models: whether $b_N$ is a property of every cascade of the class or a fitted constant of each model. The missing step is a proof, and the quantity it limits is unsized.
\end{workedexample}

\noindent The same reading of the zero-phonon to phonon-assisted ratio in \cref{sec:13.9-open} gives a different verdict. There the strongest statement is an estimate, and the missing step, the ratio of the two rates on common bins, moves the mass at which the zero-phonon line of silicon takes over from the phonon-assisted rate.
